\SourcePage{147}{150}
\section{Fall of a particle to the center}

To clarify certain features of quantum-mechanical motion, it is useful to
study a case which, although it has no immediate physical realization,
involves a particle moving in a field whose potential energy becomes infinite
at a point, the origin, according to $U(r)\approx-\beta/r^2$ with $\beta>0$.
The form of the field far from the origin will not matter. We saw in \secref{18}
that this case lies precisely between those admitting ordinary stationary
\SourcePage{148}{151}
states and those in which the particle falls to the origin.

Near the origin the Schrödinger equation is
\begin{equation}
  R''+\frac2rR'+\frac\gamma{r^2}R=0,
  \tag{35.1}
\end{equation}
where $R(r)$ is the radial wave function and
\begin{equation}
  \gamma=\frac{2m\beta}{\hbar^2}-l(l+1).
  \tag{35.2}
\end{equation}
All terms of lower order in $1/r$ have been omitted. The energy $E$ is
assumed finite, so its term has likewise been dropped.

Put $R\sim r^s$. Then $s$ satisfies
\[
  s(s+1)+\gamma=0,
\]
with roots
\begin{equation}
  s_1=-\frac12+\sqrt{\frac14-\gamma},\qquad
  s_2=-\frac12-\sqrt{\frac14-\gamma}.
  \tag{35.3}
\end{equation}

For the subsequent analysis, isolate a small region of radius $r_0$ about the
origin and replace $-\gamma/r^2$ inside it by the constant
$-\gamma/r_0^2$. After determining the wave functions in this cut-off field,
we shall examine the limit $r_0\to0$.

First suppose $\gamma<1/4$. Then $s_1$ and $s_2$ are real negative numbers,
with $s_1>s_2$. For $r>r_0$, the general small-$r$ solution is
\begin{equation}
  R=Ar^{s_1}+Br^{s_2},
  \tag{35.4}
\end{equation}
where $A$ and $B$ are constants. For $r<r_0$, the solution finite at the
origin of
\[
  R''+\frac2rR'+\frac\gamma{r_0^2}R=0
\]
is
\begin{equation}
  R=C\frac{\sin kr}{r},\qquad k=\frac{\sqrt\gamma}{r_0}.
  \tag{35.5}
\end{equation}
At $r=r_0$, both $R$ and $R'$ must be continuous. One condition may
conveniently be written as continuity of the logarithmic derivative of $rR$,
\SourcePage{149}{152}
which gives
\[
  \frac{A(s_1+1)r_0^{s_1}+B(s_2+1)r_0^{s_2}}
       {Ar_0^{s_1+1}+Br_0^{s_2+1}}=k\cot kr_0,
\]
or
\[
  \frac{A(s_1+1)r_0^{s_1}+B(s_2+1)r_0^{s_2}}
       {Ar_0^{s_1}+Br_0^{s_2}}
  =\sqrt\gamma\cot\sqrt\gamma.
\]
Solving for $B/A$ yields an expression of the form
\begin{equation}
  \frac BA=\operatorname{const}\cdot r_0^{s_1-s_2}.
  \tag{35.6}
\end{equation}

As $r_0\to0$, $B/A\to0$ because $s_1>s_2$. Of the two solutions of (35.1)
that diverge at the origin, one must therefore choose the less rapidly
divergent:
\begin{equation}
  R=A\frac1{r^{|s_1|}}.
  \tag{35.7}
\end{equation}

Now let $\gamma>1/4$. The roots are complex:
\[
  s_1=-\frac12+i\sqrt{\gamma-\frac14},\qquad s_2=s_1^*.
\]
Repeating the argument again gives (35.6), which now becomes
\begin{equation}
  \frac BA=\operatorname{const}\cdot r_0^{i\sqrt{4\gamma-1}}.
  \tag{35.8}
\end{equation}
This expression has no definite limit as $r_0\to0$, so a direct limiting
procedure is impossible. Taking (35.8) into account, a real solution has the
general form
\begin{equation}
  R=\operatorname{const}\cdot\frac1{\sqrt r}
  \cos\left(\sqrt{\gamma-\frac14}\ln\frac r{r_0}
  +\operatorname{const}\right).
  \tag{35.9}
\end{equation}
The number of zeros of this function grows without bound as $r_0$ decreases.
On the one hand, (35.9) applies at sufficiently small $r$ for every finite
particle energy $E$; on the other, the ground-state wave function must have
no zeros. We therefore conclude that
\SourcePage{150}{153}
the ground state in this field corresponds to $E=-\infty$. In every
discrete-spectrum state the particle is found mainly where $E>U$. Thus as
$E\to-\infty$ it is confined to an infinitesimal region about the origin:
the particle falls to the center.

The critical field $U_{\mathrm{cr}}$ at which fall to the center becomes
possible corresponds to $\gamma=1/4$. The smallest coefficient of $-1/r^2$
occurs for $l=0$, and hence
\begin{equation}
  U_{\mathrm{cr}}=-\frac{\hbar^2}{8mr^2}.
  \tag{35.10}
\end{equation}

Formula (35.3), with $s_1$, shows that the admissible Schrödinger solution
near a point where $U\sim1/r^2$ diverges no faster than $1/\sqrt r$ as
$r\to0$. If the field becomes infinite more slowly than $1/r^2$, then near
the origin $U(r)$ may be neglected relative to the remaining terms. The
free-motion solutions result, namely $\psi\sim r^l$ (\secref{33}). If instead the
field diverges faster than $1/r^2$, as does $-1/r^s$ with $s>2$, the wave
function near the origin is proportional to $r^{s/4-1}$; see the problem in
\secref{49}. In every such case $r\psi$ vanishes at $r=0$.

Next consider a field that at large distances decreases as
$U\approx-\beta/r^2$ but is arbitrary at small distances. Suppose first that
$\gamma<1/4$. Only finitely many negative-energy levels can then
exist\footnote{At small $r$ the field is assumed such that the particle does
not fall to the center.}. Indeed, at $E=0$ the large-distance Schrödinger
equation is (35.1), with general solution (35.4). This function has no zeros
for $r\ne0$. All zeros of the required radial wave function therefore lie at
finite distances from the origin, and their number is finite. In other
words, the ordinal number of the level $E=0$ that closes the discrete
spectrum is finite.

For $\gamma>1/4$, by contrast, the discrete spectrum contains infinitely
many negative-energy levels. At $E=0$ the large-distance wave function has
the form (35.9), with infinitely many zeros, so its ordinal number is
infinite.
\SourcePage{151}{154}
Finally, let $U=-\beta/r^2$ throughout space. If $\gamma>1/4$, the particle
falls to the center. If $\gamma<1/4$, there are no negative-energy levels at
all. The $E=0$ wave function is then of the form (35.7) everywhere and has no
zeros at finite distances; it therefore corresponds to the lowest level for
the specified $l$.
